\lecture{10}{Thu 08 Feb 2024 13:41}{Subdivision}

\begin{definition}[Affine Map]
	\label{defn:AffineMap}
	An affine simpliex in $Y$ is an affine map if $\sigma: \Delta_n \to  Y$ satisfies $\sigma\left( \sum_{i = 0}^n t_i e_i \right)  = \sum_{i = 0}^n t_i \sigma(e_i)$ for any $\sum t_i = 1$.
\end{definition}

Now, let $\text{Sin}^{\text{aff}}_n$ denote the set of affine $n$-simplices in $Y$. We may define $\text{S}^{\text{aff}}_n$. Now observe that  $d$ maps from $\text{Sin}^{\text{aff}}_n \to \text{Sin}^{\text{aff}}_{n - 1}$, so we have the  \textit{affine chain-complex} $\text{Sin}^{\text{aff}}_*(Y) $. 


 \begin{definition}[Barycenter]
	\label{defn:Barycenter}
	For $\sigma: \Delta_n \to Y$, say $\sigma = [v_0, \ldots, v_n]$. The \textit{barycenter} of $\sigma$ is average of its vertices:
	\[
		b_\sigma := \frac{1}{n + 1} \sum_{i = 0}^n v_i
	.\] 
\end{definition}

\begin{definition}[Cone Construction]
	\label{defn:ConeConstruction}
	$Y \subset \mathbb{R}^k$ is a convex subspace, $b \in Y$. We define
	\begin{align*}
		b^*: \text{Sin}^{\text{aff}}_{n - 1}(Y) &\longrightarrow \text{Sin}^{\text{aff}}_n(Y) \\
	\sigma = [v_0, \ldots, v_{n - 1}] &\longmapsto b^*\sigma = [b, v_0, \ldots, v_{n - 1}]
.\end{align*}

Extend linearly to  $S^{\text{aff}}_{n - 1} \to S^{\text{aff}}_{n }$.
\end{definition}

\begin{definition}[Subdivision]
	\label{defn:Subdivision}
	We define the subdivision operator $S':S^{\text{aff}}_{n } \to S^{\text{aff}}_{n }$ such that $S' = \text{Id}$ for $n \le 0$, and for $n > 0$,
	\[
		S'(\sigma) = b_\sigma * S'(d \sigma)
	.\] 
\end{definition}

\begin{example}
	For $n = 1$, $\sigma = [v_0 v_1]$, $b_\sigma = \frac{1}{2}(v_0 + v_1)$. Now
	\[
		S'(\sigma) = b_\sigma * S'([v_1] - [v_0]) = b_\sigma* ([v_1] - [v_0]) = [b_\sigma v_1] - [b_\sigma v_0]
	.\] 
\end{example}


\begin{remark}
	For each sequence $f_0 \supset f_1\supset\ldots\supset f_n$, where $f_0 = \sigma$, $f_{i + 1}$ is a face of $f_i$, $f_n$ is a vertex of $\sigma$. 

	Let $b_i$ be barycenter of $f_i$. Then $[b_0, b_1, \ldots, b_n] \in  S'(\sigma)$. From this argument we see that the number of $n$-simplices in $S'(\sigma)$ is $(n + 1)!$.
\end{remark}


\begin{lemma}
	\label{lem:SubdivisionChainMap}
	$S'$ is a chain map, i.e. the following diagram commutes.
% https://quiver.local:8080/#q=WzAsNCxbMCwwLCJTXntcXHRleHR7YWZmfX1fe24gfSAiXSxbMiwwLCJTXntcXHRleHR7YWZmfX1fe24gfSAiXSxbMCwyLCJTXntcXHRleHR7YWZmfX1fe24gLSAxfSAiXSxbMiwyLCJTXntcXHRleHR7YWZmfX1fe24gLSAxfSAiXSxbMCwyLCJkIl0sWzEsMywiZCJdLFswLDEsIlxcJCIsMV0sWzIsMywiXFwkIiwxXV0=
\[\begin{tikzcd}
	{S^{\text{aff}}_{n } } && {S^{\text{aff}}_{n } } \\
	\\
	{S^{\text{aff}}_{n - 1} } && {S^{\text{aff}}_{n - 1} }
	\arrow["d", from=1-1, to=3-1]
	\arrow["d", from=1-3, to=3-3]
	\arrow["{\$}"{description}, from=1-1, to=1-3]
	\arrow["{\$}"{description}, from=3-1, to=3-3]
\end{tikzcd}\]
\end{lemma}
\begin{proof}
	Use induction. For $n = 0$, they are all identity operators. Now assume $d \$ = \$ d$ for $n = k - 1, k \ge 1$. Now for $n = k$, $\sigma$ be an affine $n$-simplex.
	\begin{align*}
		d \$ (\sigma) &= d(b_\sigma * \$ d \sigma) \\
		&= \$ d \sigma - b_\sigma * d(\$ d \sigma) \\
		&= \$ d \sigma
	.\end{align*}
	For last equality, note that IH gives $d \$ (d \sigma) = \$ d (d \sigma) = \$ (d d) \sigma = 0$.

\end{proof}

